% Source: https://artofproblemsolving.com/community/c3h513600_mathcounts_formulas
% Stable identity: aops-topic-513600-q01
\textbf{The formula.} Let $I$ be the incentre of triangle $ABC$, with sides $a=BC$, $b=CA$, $c=AB$. The incircle touches each side, so the distance from $I$ to each side is $r$. The segments $IA,IB,IC$ cut the triangle into triangles $IBC$, $ICA$, $IAB$ with bases $a,b,c$ and common height $r$:
\[
K=\dfrac{1}{2}ar+\dfrac{1}{2}br+\dfrac{1}{2}cr=r\cdot\dfrac{a+b+c}{2}=rs,\qquad\text{so}\qquad r=\dfrac {K}{s} .\qquad\blacksquare
\]

\textbf{The equilateral triangle.} With side $4$, $K=\dfrac{\sqrt3}{4}\cdot4^2=4\sqrt3$ and $s=6$:

\begin{asy}
usepackage("amsmath");
size(4.5cm);
defaultpen(fontsize(10pt));
real a=4; pair B=(0,0),C=(a,0),A=(a/2,a*sqrt(3)/2),I=(a/2,a*sqrt(3)/6),M=(a/2,0);
draw(A--B--C--cycle,linewidth(0.9));
draw(circle(I,a*sqrt(3)/6),blue);
draw(I--M,red);real s=0.18; draw(M+(s,0)--M+(s,s)--M+(0,s));
dot(I);
label("$A$",A,N);label("$B$",B,SW);label("$C$",C,SE);label("$I$",I,N);
label("$r$",(I+M)/2,W,red);label("$4$",(A+C)/2,NE);
label("\mbox{}",(min(currentpicture,user=true)+max(currentpicture,user=true))/2,black); // keeps the text after the figure black when LaTeX typesets the labels inline
\end{asy}

\[
r=\dfrac{4\sqrt3}{6}=\boxed{\dfrac{2\sqrt3}{3}}.
\]
(Check: in an equilateral triangle the incentre is the centroid, so $r$ is one third of the height $2\sqrt3$.)
